\documentclass[spanish,letterpaper,oneside]{article}

\def\documenttitle {Primer Momento: Análisis del Aprendizaje Profundo (Quiroga y Lara, 2022)}
\def\documentsubtitle {Fundamentos para la enseñanza y el aprendizaje I · Subproducto 6A (Bloque III)}
\def\documentsubject {Actividad 6 · Primer Momento}
\def\documentauthor {Mtro. Martín Jonathan de la Cruz Muñoz}
\def\coursename {Fundamentos para la enseñanza y el aprendizaje I}
\def\coursecode {MAEM-SD26}

\def\universityname {Instituto de Investigación, Innovación y Estudios de Posgrado para la Educación del Estado de Nuevo León}
\def\universityfaculty {Maestría en Aprendizaje y Enseñanza de las Matemáticas}
\def\universitydepartment {Fundamentos para la enseñanza y el aprendizaje I}
\def\universitydepartmentimage {img/iiiepe/marca-iiiepe.png}
\def\universitydepartmentimagecfg {height=1.75cm}
\def\universitylocation {Monterrey, Nuevo León}
\def\studentfullname {Mtro. Martín Jonathan de la Cruz Muñoz}
\def\studentid {00874}
\def\studentemail {martin.delacruz@campus.iiiepe.edu.mx}
\def\programperiod {2do periodo, septiembre--diciembre de 2026 (SD26)}
\def\programmodality {Modalidad presencial}
\def\programshort {MAEM}

\def\authortable {
  \begin{tabular}{ll}
    \textbf{Alumno:} & \studentfullname \\
    Matrícula: & \studentid \\
    Correo institucional: & \studentemail \\
    Unidad: & \coursename \\
    Programa: & \programshort \\
    Periodo: & \programperiod \\
    & \\
    \multicolumn{2}{l}{Fecha de realización: 19 de noviembre de 2026} \\
    \multicolumn{2}{l}{\universitylocation}
  \end{tabular}
}

\input{template}

\usepackage{helvet}
\renewcommand{\familydefault}{\sfdefault}
\usepackage{array}
\usepackage{longtable}
\usepackage{pdflscape}
\usepackage{tikz}
\usepackage{pgfplots}
\pgfplotsset{compat=1.18}
\usepackage[most]{tcolorbox}
\usetikzlibrary{arrows.meta,positioning,calc,fit,backgrounds,shapes.geometric}

\definecolor{iiiepeNavy}{HTML}{163B5C}
\definecolor{iiiepeBlue}{HTML}{1F6F9F}
\definecolor{iiiepeAqua}{HTML}{20B7A8}
\definecolor{iiiepeOrange}{HTML}{F28C28}
\definecolor{iiiepePaper}{HTML}{F6F8F7}
\definecolor{iiiepeInk}{HTML}{1E2A32}
\definecolor{iiiepePale}{HTML}{EEF5F8}
\definecolor{iiiepeMint}{HTML}{E7F6F4}
\definecolor{iiiepeWarm}{HTML}{FFF3E6}
\def\portraittitlecolor {iiiepeNavy}
\setcitestyle{authoryear,open={(},close={)}}
\makeatletter
\renewcommand{\@biblabel}[1]{}
\makeatother
\setlength{\bibhang}{1.27cm}
\setlength{\bibsep}{6pt}
\AtBeginDocument{\renewcommand{\refname}{Referencias}}

\newcommand{\insertcoverwatermark}{%
  \AddToShipoutPictureBG*{%
    \begin{tikzpicture}[remember picture,overlay]
      \node[opacity=0.04,inner sep=0pt] at (current page.center) {%
        \includegraphics[width=0.65\paperwidth]{img/iiiepe/marca-iiiepe.png}%
      };
    \end{tikzpicture}%
  }%
}

\setlength{\parskip}{5pt}\setlength{\parindent}{0pt}
\begin{document}\insertcoverwatermark\templatePortrait\templatePagecfg\templateFinalcfg\fancyfoot{}\fancyfoot[L]{\scriptsize IIIEPE · MAEM · SD26}\fancyfoot[C]{\scriptsize\thepage}\fancyfoot[R]{\scriptsize Subproducto 6A · Primer Momento}\setstretch{1.3}\raggedbottom
\section{Primer momento: análisis teórico de la lectura base}
Como primer momento de fundamentación teórica, este estudio examina la investigación de \citet{quiroga2022aprendizaje}, \textit{El aprendizaje profundo como herramienta para cambio en la visión de aprendizaje de una cultura escolar}.

Mi postura es que transformar las clases de matemáticas en bachillerato técnico e ingeniería resulta inviable si se modifican únicamente formatos externos sin revisar la concepción profunda del aprendizaje, los vínculos de poder en el aula y las creencias docentes sobre la cultura escolar \citep{quiroga2022aprendizaje,sawyer2008optimizing}.

\section{Cinco ejes críticos analizados}
A partir de la revisión sistemática de \citet{quiroga2022aprendizaje} en diálogo con la tradición cognitiva y de nuevas pedagogías \citep{marton1976qualitative,fullan2014rich}, desgloso cinco pilares conceptuales fundamentales:

\begin{enumerate}
  \item \textbf{Diferencias entre aprendizaje superficial y profundo:} Recuperando la distinción seminal de \citet{marton1976qualitative}, el aprendizaje superficial busca memorizar reglas, pasos y definiciones de manera fragmentada para acreditar una evaluación puntual, produciendo saberes rígidos e intransferibles. En contraste, el aprendizaje profundo integra conocimiento previo con nuevas ideas, exige comprender razones subyacentes y habilita aplicar modelos matemáticos frente a situaciones nuevas \citep{sawyer2008optimizing}.
  \item \textbf{Integración de las dimensiones cognitiva, interpersonal e intrapersonal:} Para \citet{quiroga2022aprendizaje}, la comprensión duradera excede el plano meramente intelectual. Exige articular procesos \textit{cognitivos} (modelación, análisis, solución de problemas), capacidades \textit{interpersonales} (comunicación argumentada, escucha cooperativa y negociación) y disposiciones \textit{intrapersonales} (metacognición, perseverancia ante el bloqueo y mentalidad de crecimiento) \citep{schunk2012teorias}.
  \item \textbf{Transformación en los roles docente y estudiantil:} Retomando el marco de \citet{fullan2014rich}, el profesorado abandona el papel de emisor único para actuar como diseñador de ambientes retadores, investigador de obstáculos estudiantiles y co-aprendiz que construye relaciones de asociación auténtica. Paralelamente, el alumnado asume protagonismo en formular preguntas, validar pruebas y responsabilizarse colectivamente de los acuerdos de equipo.
  \item \textbf{Características de las nuevas pedagogías y culturas escolares:} Se caracterizan por problemas del mundo real en lugar de ejercicios artificiales, uso reflexivo de tecnología para potenciar indagación en vez de sustituir pensamiento, asociaciones explícitas docente--estudiante y liderazgo orientado por metas comunes \citep{rivas2017innovaciones}.
  \item \textbf{Creencias pedagógicas como núcleo del cambio cultural:} Toda innovación fracasa si persisten creencias implícitas según las cuales «enseñar bien significa cubrir muchos ejercicios idénticos en silencio» o «quien no trae base algebraica fracasa inevitablemente». Por consiguiente, modificar la cultura escolar exige cuestionar primero las creencias que organizan nuestra práctica cotidiana \citep{quiroga2022aprendizaje}.
\end{enumerate}

Además, el cuadro siguiente sintetiza cómo estos cinco ejes interpelan directamente nuestras prácticas de enseñanza matemática en nivel medio superior técnico y superior:

{\small\setstretch{1.15}\renewcommand{\arraystretch}{1.3}
\begin{longtable}{|>{\raggedright\arraybackslash}p{0.24\linewidth}|>{\raggedright\arraybackslash}p{0.35\linewidth}|>{\raggedright\arraybackslash}p{0.35\linewidth}|}\hline
\rowcolor{iiiepePale} \textbf{Eje analítico} & \textbf{Inercia escolar habitual} & \textbf{Horizonte del aprendizaje profundo} \\\hline
\endfirsthead
\hline
\rowcolor{iiiepePale} \textbf{Eje analítico} & \textbf{Inercia escolar habitual} & \textbf{Horizonte del aprendizaje profundo} \\\hline
\endhead
Naturaleza del saber & Procedimientos aislados por capítulo. & Red de conceptos transferibles en contexto \citep{sawyer2008optimizing}. \\\hline
Dimensiones activadas & Foco exclusivo en ejecución cognitiva individual. & Integración cognitiva, interpersonal e intrapersonal \citep{quiroga2022aprendizaje}. \\\hline
Relación pedagógica & Transmisión y obediencia directiva. & Asociación docente--estudiante y autonomía guiada \citep{fullan2014rich}. \\\hline
Sentido de evaluación & Calificación sumativa final sin retorno. & Devolución formativa y reflexión metacognitiva \citep{eef2021feedback}. \\\hline
\end{longtable}}

\clearpage
\section{Conclusiones}
En conclusión, considero que el análisis teórico desarrollado en este Primer momento aporta la claridad conceptual indispensable antes de diagnosticar la propia práctica y diseñar innovaciones concretas. Sin embargo, dado que ninguna categoría teórica mejora por sí sola el aprendizaje matemático en clase, además es preciso vincular la dimensión cognitiva con el diálogo interpersonal; es decir, por ello sostengo que asumir el aprendizaje profundo como transformación cultural nos protege de aplicar estrategias aisladas sin cambiar el sentido ético y cognitivo del aula.

\clearpage
\begin{thebibliography}{99}
\bibitem[Collin y Quigley(2021)]{eef2021feedback} Collin, J., y Quigley, A. (2021). \textit{Teacher feedback to improve pupil learning}. Education Endowment Foundation.
\bibitem[Fullan y Langworthy(2014)]{fullan2014rich} Fullan, M., y Langworthy, M. (2014). \textit{A rich seam: How new pedagogies find deep learning}. Pearson.
\bibitem[Marton y Säljö(1976)]{marton1976qualitative} Marton, F., y Säljö, R. (1976). On qualitative differences in learning: I. Outcome and process. \textit{British Journal of Educational Psychology}, 46(1), 4--11.
\bibitem[Quiroga Aguilar y Lara Órdenes(2022)]{quiroga2022aprendizaje} Quiroga Aguilar, L., y Lara Órdenes, E. (2022). El aprendizaje profundo como herramienta para cambio en la visión de aprendizaje de una cultura escolar. \textit{Revista Educación las Américas}, 12(1).
\bibitem[Rivas et~al.(2017)]{rivas2017innovaciones} Rivas, A., André, F., y Delgado, L. E. (Eds.). (2017). \textit{50 innovaciones educativas para escuelas}. CIPPEC / Santillana.
\bibitem[Sawyer(2008)]{sawyer2008optimizing} Sawyer, R. K. (2008). Optimizing learning: Implications of learning sciences research. OECD Publishing.
\bibitem[Schunk(2012)]{schunk2012teorias} Schunk, D. H. (2012). \textit{Teorías del aprendizaje} (6.ª ed.). Pearson.
\end{thebibliography}

\end{document}
